\section{Domain and Theoretical Foundations}\label{sec:domain-foundations}

\subsection*{The security-operations domain}

The domain this dissertation addresses is offensive and defensive security
operations: the systematic discovery, assessment and (where authorised)
remediation of weaknesses in networks and applications. The operator's work is
organised around a small set of recurring intents---discovering what is
exposed, enumerating services, identifying known vulnerabilities, and checking
whether specific weaknesses apply to particular targets. These intents map onto
specialised tooling: network discovery and scanning (of which Nmap is the
canonical example \parencite{lyon2009nmap}), and exploitation framework
drivers (of which Metasploit is the canonical example \parencite{metasploit}).
Two properties of this domain matter for the work in this dissertation. First,
each intent carries a \emph{target} and a set of \emph{parameters}, so the gap
between prose and machine dispatch is structural rather than cosmetic. Second,
every action is bounded by authorisation: a scan may only touch targets and
techniques within an explicitly declared engagement scope, and some intents are
reserved for roles with the required mandate. The translation layer evaluated
in this dissertation is exactly the component that must preserve both the
structure and the authorisation of operator intent.

\subsection*{Intent recognition and natural-language understanding}

At its core, intent recognition is a natural-language understanding (NLU)
problem: given an utterance, determine the communicative intent and extract the
arguments that parameterise it. In a classical NLU system this is split into an
intent-classification step, which assigns the utterance to one of a fixed set
of intent classes, and a slot-filling (or entity-extraction) step, which
recovers typed argument values such as hosts, networks and ports. The intent
contract introduced in Section~\ref{sec:individual-component-context} is the
machine-consumable record of these two steps: a resolved intent type, a target
and recognised options, together with any validation findings. For this record
to be useful to downstream automation it must be expressed in a
schema-bounded, machine-parseable syntax, which is why the IRE emits a
structured contract rather than prose.

Representational infrastructure underpins the more advanced parts of the
design. Sentence embeddings map natural-language strings into a vector space in
which semantically similar strings are close under a distance such as cosine
similarity; Sentence-BERT is the representative technique \parencite{reimers2019sbert}.
These embeddings provide the two similarity signals used in this work: the
agreement signal by which multiple model outputs are judged to express the same
meaning, and the cache signal by which a repeated intent is recognised without
re-invoking the model.

\subsection*{Language models as stochastic components}

Instruction-tuned language models are the generative core of the IRE. Their
defining property for this dissertation is that inference is a sampling
process: given a prompt, the model defines a probability distribution over
output sequences, and a response is drawn from that distribution. Sampling
temperature controls how peaked the distribution is, so the same prompt can
yield different outputs across calls. A single response is therefore a
\emph{sample} of the model's belief, not a deterministic verdict. This
observation is the theoretical basis of the design's treatment of the model as
an unreliable proposition generator whose output must be checked rather than
trusted. It also underlies consistency-based confidence: the
self-consistency decoding strategy samples several responses and strengthens
confidence in answers that agree with one another \parencite{wang2023selfconsistency},
applying a voting intuition that originates in chain-of-thought reasoning and
transfers naturally to classification agreement.

\subsection*{Hallucination as a failure class}

A second property of language models is that their outputs can be fabricated
relative to the true state of the world. The hallucinations relevant here are
the \emph{extrinsic} form in the taxonomy of Ji \emph{et al.}: statements that
contradict or are unverifiable against an authoritative source, rather than
mere inconsistency with an input context \parencite{ji2023hallucination}.
Fabricated security identifiers are a concrete instance: a model may emit a
CVE identifier, an IP address or a port that does not exist or is not the one
the operator intended. Because a downstream tool will act on those values, the
practical requirement is not simply to detect fabrication in general but to
detect it in tightly constrained value spaces (identifier formats, address
syntaxes, port ranges), where deterministic checking is feasible.

\subsection*{Agreement and uncertainty as fabrication signals}

Detection without an external verifier relies on the statistical structure of
the model's own output. Sampling-based methods compute a consistency score from
a small set of generated responses---SelfCheckGPT scores sampled responses for
mutual consistency \parencite{manakul2023selfcheckgpt}---while information-theoretic
approaches frame the uncertainty of free-form answers through semantic
entropy over meaning-equivalence clusters \parencite{kuhn2023semantic,
farquhar2024semantic}. These methods provide the theoretical grounding for the
self-consistency gate in the IRE design. Their common cost, however, is
repeated generation: every extra sample multiplies inference time, which
collides with the CPU-only, single-user latency budget adopted in this
dissertation. That tension---between the statistical value of sampling and the
hard cost of each sample---is a central theme returned to in the synthesis
(Section~\ref{sec:gap-synthesis}).

\subsection*{Authorisation and network boundaries}

Finally, two security foundations bound what the validated contract is allowed
to assert. The first is role-based access control (RBAC), under which a user's
authorisation to perform an action is determined by the roles assigned to that
user and an access matrix that grants or denies the intent for each role; the
IRE's four-role model operationalises this standard pattern. The second is
network-level scope: private address allocation is codified by RFC~1918
\parencite{rfc1918}, which defines the address ranges reserved for private
internets and provides the technical basis for distinguishing an internal
target from an external one. Together with the operator-declared engagement
scope, these foundations define the predicates that deterministic validation
can and cannot decide: syntax, format and authorisation are decidable
propositions; only meaning is delegated to the model. This division of
responsibility is the theoretical spine of the solution designed in
Chapter~\ref{ch:product-design}.

\TODO{If the assessment rubric requires canonical textbook references for NLU
fundamentals and for role-based access control, add formal citations for the
intent-classification/slot-filling formulation and for the RBAC reference model;
none are currently in the reference database.}